\documentclass[aspectratio=169,11pt]{beamer}
\usepackage{remix}

\begin{document}

%% ---------- TITLE ----------
\remixtitleframe
  {Advanced Topics in Applied Causal Inference}
  {A test of the Remix style}
  {Scott Cunningham, Baylor University}
  {University of Glasgow \textperiodcentered\ May 18--22, 2026}

%% ---------- SECTION BREAK ----------
\remixsection{1. DiD Origins}

%% ---------- CONTENT FRAME ----------
\begin{frame}{The four averages and three subtractions}{Ashenfelter's move at Princeton}
\begin{itemize}
  \item DiD is not a regression. It's an \highlight{idea}.
  \item Orley Ashenfelter at the Industrial Relations Section: applied research questions, communicate to bureaucrats, get the answer right with high-quality data.
  \item His move: don't emphasize the regression --- emphasize the simplicity of \highlight{four averages and three subtractions}.
  \item The 2$\times$2 table: treated/control by pre/post; four cells; four averages.
\end{itemize}

\vspace{0.5em}

\begin{block}{The DiD formula}
\[
\widehat{\text{DiD}} = \underbrace{(Y_{11} - Y_{10})}_{\text{treated change}} - \underbrace{(Y_{01} - Y_{00})}_{\text{control change}}
\]
\end{block}
\end{frame}

%% ---------- BLOCK SHOWCASE ----------
\begin{frame}{Block styles}{Three flavors for three uses}
\begin{block}{Standard block}
Body text uses slate. Title is charcoal on a light background. For neutral commentary or definitions.
\end{block}

\begin{alertblock}{Alert block}
For \highlight{key insights} or warnings. The accent color tints the background subtly. Use sparingly --- restraint matters.
\end{alertblock}

\begin{exampleblock}{Example block}
For worked examples or empirical illustrations. Ocean-blue tint distinguishes it from the alert.
\end{exampleblock}
\end{frame}

%% ---------- PULL QUOTE ----------
\begin{frame}{Voice of the field}
\pullquote{If you can convince a bureaucrat in five minutes, you have done something. Anything else is academic.}{Orley Ashenfelter (paraphrased, Scott's creative license)}

\vspace{1em}
The point is communication. The 2$\times$2 doesn't need a regression to tell its story.
\end{frame}

%% ---------- SECTION BREAK 2 ----------
\remixsection{2. The Three Regressions}

%% ---------- TWO-COLUMN ----------
\begin{frame}{Equivalence on castle.dta}{2005--2006 subset, one DiD, three ways}
\twocol{%
\textbf{Three specifications:}
\begin{enumerate}
  \item OLS with interactions \\ \texttt{reg l\_homicide post\#\#treat}
  \item TWFE \\ \texttt{xtreg ..., fe vce(cluster sid)}
  \item Long difference \\ \texttt{reshape wide; reg diff treat}
\end{enumerate}
}{%
\textbf{All three return the same number} to machine epsilon.

\vspace{0.5em}
This equivalence is what \highlight{made it possible} to extend DiD to covariates, staggered, and continuous treatment using OLS infrastructure.

\vspace{0.5em}
\meta{But the equivalence breaks once you add complexity. Day 2.}
}
\end{frame}

\end{document}
